\scanpage{24}
\mspage{18}

\[
 C,TC\qquad(C\subset\mathrm{Cof}_W,\ TC\subset\deleted{\unsure{W^{\mathrm{univ}}\cap}}C\cap W)
\]
donnant naissance à
\[
 F=(TC)_*,\quad TF=C_*,
\]
par
\[
 C,\ C\cap W\quad(\supset TC)
\]
ce qui ne change pas le couple $(C,TF)$,
mais bien le couple $(TC,F)$, remplacé
par
\[
 \begin{aligned}
 (TC)' &={} C\cap W,& F'&=(TC')_*\\[-2pt]
 \cup &\deleted{\unsure{\cap|?|\,\uparrow}}\hspace{1.5em}\parallel &&\cap\\[-2pt]
 TC &\quad C\cap W^{\mathrm{univ}}&&F
 \end{aligned}
\]
% slanted
\margin{NB $TC'$ sera encore clos à gauche, comme intersection de deux parties closes à gauche.}
\unsure{--} \unsure{renforcé|remplacé} $F$ par un $F'$ plus petit).

Le seul problème, c'est si $\deleted{\unsure{(C,TF)}}$
$(TC',F')$ satisfait encore à la condition
de factorisation de Quillen. Ce sont des
parties closes (à gauche pour $TC'$, à dr.
pour $F'$) ``duales'' l'une de l'autre
\[
 TC'={F'}^{*},\quad F'=TC'_*,
\]
mais comme on ne sait pas \deleted{\unsure{d'avance|d'avan}}\inserted{si}
$TC'$ est \deleted{\unsure{engendré}} de la forme $\widetilde{TC'_0}$, avec
$TC'_0$ petite, on n'est pas sûr que la \deleted{\unsure{th.}}
conclusion du th.~1 s'applique. \deleted{\unsure{Mais}} \inserted{Mais}
si on suppose dans a$'$) que $C$ est
une partie \uline{accessible} de $\mathrm{Fl}\,M$ (i.e., étant
déjà close à g, que $C=\widetilde{C_0}$, $C_0$ petite),
% archive sheet number: 23
